% Part: first-order-logic
% Chapter: model-theory
% Section: isomorphism

\documentclass[../../../include/open-logic-section]{subfiles}

\begin{document}

\olfileid{mod}{bas}{iso}
\olsection{Struktur-Struktur Isomorfis}

!!{structure}s orde kapisan bisa padha
ing salah siji saka rong cara.
Cara kapisan yaiku padha ndadekake
!!{sentence}s kang padha dadi bener.
!!{structure}s kaya mangkono
diarani \emph{ekuivalen elementer}.
Nanging loro struktur bisa beda
banget lan isih ndadekake
!!{sentence}s kang padha
dadi bener---upamane, sing siji
bisa !!{enumerable}, dene
sijine ora. Iki amarga akeh
sipat !!a{structure} kang
ora bisa dinyatakake ing
basa orde kapisan, bisa
amarga basane kurang sugih,
utawa amarga watesan
dhasar logika orde kapisan,
kayata teorema
L\"owenheim--Skolem.
Mula ana cara kapindho
kang luwih ketat kanggo
ngarani rong !!{structure}s
padha: loro-lorone
sajatiné padha, mung
objek-objek kang nyusune
beda, dudu sipat
strukturalé. Pangertèn
iki bisa diprecisake
kanthi \emph{isomorfisme}.

\begin{defn}
\ollabel{defn:elem-equiv}
Yen ana loro !!{structure}s
$\Struct{M}$ lan $\Struct M'$
kanggo !!{language}~$\Lang{L}$
kang padha, $\Struct{M}$
diarani \emph{ekuivalen elementer
  karo} $\Struct M'$,
ditulis $\Struct{M} \equiv \Struct M'$,
yen lan mung yen kanggo
saben !!{sentence}~$!A$
saka~$\Lang{L}$,
$\Sat{M}{!A}$ yen lan
mung yen $\Sat{M'}{!A}$.
\end{defn}

\begin{defn}
\ollabel{defn:isomorphism}
Yen ana loro !!{structure}s
$\Struct{M}$ lan
$\Struct M'$ kanggo
!!{language}~$\Lang L$
kang padha, $\Struct{M}$
diarani \emph{isomorfis
karo}~$\Struct M'$,
ditulis $\Struct{M}
\simeq \Struct M'$,
yen lan mung yen ana
fungsi $h \colon
\Domain{M} \to \Domain{M'}$
kanthi sipat-sipat iki:
\begin{enumerate}
\item $h$ iku !!{injective}:
  yen $h(x) = h(y)$
  mula $x = y$;
\item $h$ iku !!{surjective}:
  kanggo saben $y \in \Domain{M'}$
  ana $x \in \Domain{M}$
  saéngga $h(x) = y$;
\item \ollabel{defn:iso-const}kanggo saben
  !!{constant} $c$:
  $h(\Assign{c}{M}) = \Assign{c}{M'}$;
\item \ollabel{defn:iso-pred}kanggo saben
  $n$-panggonan !!{predicate}~$P$:
  \[
  \tuple{a_1, \dots, a_n}\in \Assign{P}{M} \quad\text{yen lan mung yen}\quad
  \tuple{h(a_1), \dots, h(a_n)} \in \Assign{P}{M'};
  \]
\item \ollabel{defn:iso-func}kanggo saben
  $n$-panggonan !!{function} $f$:
  \[
  h(\Assign{f}{M}(a_1, \dots, a_n)) =
  \Assign{f}{M'}(h(a_1), \dots, h(a_n)).
  \]
\end{enumerate}
\end{defn}

\begin{thm}
\ollabel{thm:isom}
Yen $\Struct{M} \iso \Struct M'$
mula $\Struct{M} \elemequiv
\Struct{M'}$.
\end{thm}

\begin{proof}
Upamakna $h$ iku isomorfisme
saka $\Struct{M}$ menyang
$\Struct M'$. Kanggo
sembarang pangancikan~$s$,
$h \circ s$ iku komposisi
$h$ lan $s$, yaiku
pangancikan ing
$\Struct{M'}$ kang nyukupi
$(h \circ s)(x) = h(s(x))$.
Kanthi induksi tumrap
$t$ lan $!A$, kita
bisa mbuktekake pratelan
kang luwih kuwat iki:
\begin{enumerate}
  \item[a.] $h(\Value{t}{M}[s]) = \Value{t}{M'}[h\circ s]$.
  \item[b.] $\Sat{M}{!A}[s]$ yen lan mung yen
    $\Sat{M'}{!A}[h \circ s]$.
\end{enumerate}
Pratelan kapisan dibuktekake
kanthi induksi tumrap
kerumitan~$t$.
\begin{enumerate}
\item Yen $t \ident c$, mula
  $\Value{c}{M}[s] = \Assign{c}{M}$
  lan $\Value{c}{M'}[h \circ s] =
  \Assign{c}{M'}$. Dadi,
  $h(\Value{t}{M}[s]) =
  h(\Assign{c}{M}) = \Assign{c}{M'}$
  (miturut \olref{defn:iso-const}
  saka \olref{defn:isomorphism})
  $= \Value{t}{M'}[h \circ s]$.
\item Yen $t \ident x$, mula
  $\Value{x}{M}[s] = s(x)$
  lan $\Value{x}{M'}[h \circ s] =
  h(s(x))$. Dadi,
  $h(\Value{x}{M}[s]) =
  h(s(x)) = \Value{x}{M'}[h \circ s]$.
\item Yen $t \ident f(t_1, \dots, t_n)$,
  mula
  \begin{align*}
    \Value{t}{M}[s] & = \Assign{f}{M}(\Value{t_1}{M}[s], \dots, \Value{t_n}{M}[s]) \quad\text{lan}\\
  \Value{t}{M'}[h \circ s] & = \Assign{f}{M'}(\Value{t_1}{M'}[h \circ
    s], \dots, \Value{t_n}{M'}[h \circ s]).
  \end{align*}
  Cathetan koreksi sumber OLPL-087: nilai term ing
  struktur kapindho nggunakake interpretasi fungsi
  ing struktur kapindho, dudu ing struktur kapisan.
  Hipotesis induksine yaiku kanggo saben $i$,
  $h(\Value{t_i}{M}[s])
  = \Value{t_i}{M'}[h\circ s]$.
  Dadi,
  \begin{align}
    h(\Value{t}{M}[s])
    & = h(\Assign{f}{M}(\Value{t_1}{M}[s], \dots, \Value{t_n}{M}[s])) \notag\\
    & = \Assign{f}{M'}(h(\Value{t_1}{M}[s]), \dots,
    h(\Value{t_n}{M}[s])) \ollabel{iso-1}\\
    & = \Assign{f}{M'}(\Value{t_1}{M'}[h \circ s], \dots,
    \Value{t_n}{M'}[h \circ s]) \ollabel{iso-2}\\
    & = \Value{t}{M'}[h\circ s] \notag
  \end{align}
  Cathetan koreksi sumber OLPL-088: baris kapisan
  rantai persamaan ing ndhuwur nutup uga kurung
  njaba kanggo aplikasi fungsi isomorfisme.
  Ing kene, \olref{iso-1}
  ndherek saka \olref{defn:iso-func}
  ing \olref{defn:isomorphism},
  lan \olref{iso-2} saka
  hipotesis induksi.
\end{enumerate}
Bagean (b) ditinggal minangka gladhen.

Yen $!A$ iku ukara,
pangancikan~$s$ lan
$h \circ s$ ora mengaruhi
kasunyatane, mula
$\Sat{M}{!A}$ yen lan
mung yen $\Sat{M'}{!A}$.
\end{proof}

\begin{prob}
Rampungna bukti bagean (b)
saka \olref[mod][bas][iso]{thm:isom}
kanthi rinci. Tuduhna
ing ngendi saben siji
saka limang sipat
isomorfisme ing
\olref[mod][bas][iso]{defn:isomorphism}
digunakake.
\end{prob}

\begin{defn}
\emph{Automorfisme} saka
struktur $\Struct{M}$
iku isomorfisme saka
$\Struct{M}$ menyang
awake dhewe.
\end{defn}

\begin{prob}
Tuduhna yen kanggo
sembarang struktur
$\Struct{M}$, yen $X$
iku subhimpunan saka
$\Struct{M}$ kang
bisa ditetepake,
lan $h$ iku automorfisme
saka $\Struct{M}$,
mula $X
= \Setabs{h(x)}{x \in X}$
(yaiku, $X$ tetep
miturut $h$).
\end{prob}

\end{document}
